% ===============================================================
% 技术方案
% ===============================================================
\section{技术方案}

% ---------------------------------------------------------------
\begin{frame}{\sa{}}
\begin{columns}[c,onlytextwidth]
\begin{column}{0.71\textwidth}
\centering
\includegraphics[width=\linewidth,height=0.95\textheight,keepaspectratio]{arch.pdf}
\end{column}
\begin{column}{0.27\textwidth}
\begin{tcolorbox}[rustcard,title={宿主机},fontupper=\footnotesize]
  构建时选择性插桩并校验边界；\code{syz-manager} 负责调度、panic 解析与 VM 生命周期
\end{tcolorbox}
\vspace{1mm}
\begin{tcolorbox}[tcbcard,title={目标机},fontupper=\footnotesize]
  插桩点回调 \code{ostd::kcov}，PC 写入\textbf{共享 VMO}，执行器零拷贝读取
\end{tcolorbox}
\vspace{1mm}
\begin{tcolorbox}[syzcard,title={工作流程},fontupper=\footnotesize]
  打补丁 $\to$ 插桩构建 $\to$ 镜像 $\to$ 模糊测试 $\to$ 分诊复现 $\to$ 统计评估
\end{tcolorbox}
\end{column}
\end{columns}
\end{frame}

% ---------------------------------------------------------------
\begin{frame}{\sa{}}
\renewcommand{\arraystretch}{1.6}
\begin{tabularx}{\textwidth}{@{}>{\bfseries}p{2.6cm}>{\raggedright\arraybackslash}X@{}}
  \toprule
  \normalfont\muted{挑战} & \muted{\sa{} 的方案} \\
  \midrule
  \textcolor{SARust}{插桩互相干扰}
    & \hl{插桩边界 = TCB 边界}：只插桩 \code{aster\_core}，禁用跨 crate LTO；
      只在任务上下文中以非重入方式记录 PC \\
  \textcolor{SARust}{观测接口缺失}
    & 与 Linux 兼容的 \code{/dev/kcov}（共享 VMO）和 \code{/proc/kallsyms}；
      执行器与符号解析自动回退 \\
  \textcolor{SASyz}{panic 无法解析}
    & 识别 \code{Uncaught panic:}，提取原因与源码位置；截断路径、归一化行号后去重 \\
  \textcolor{SASyz}{重启开销高}
    & \hl{快照库 + LinUCB 上下文多臂老虎机}：在继续、检查点、软重启、硬重启之间自适应决策 \\
  \bottomrule
\end{tabularx}
\end{frame}
